% 정답표 (A4) — make_tex.py가 answers.md를 읽어 \answersheet 내용을 채운다.
\documentclass{article}
\usepackage[a4paper,margin=25mm]{geometry}
\usepackage{fontspec}
\usepackage{kotex}
\usepackage{amsmath,array}
\pagestyle{empty}
\setmainfont{Latin Modern Roman}
\setmainhangulfont{NanumMyeongjo}[BoldFont={NanumMyeongjo Bold}]
\newfontfamily\gothicL{NanumGothic}[BoldFont={NanumGothic ExtraBold}]
\newhangulfontfamily\gothicH{NanumGothic}[BoldFont={NanumGothic ExtraBold}]
\newcommand\gothic{\gothicL\gothicH}
\newfontfamily\circfont{NanumMyeongjo}
\newcommand\answersheet[2]{%
  \begin{center}
  {\gothic\bfseries\fontsize{20bp}{26bp}\selectfont 수학 영역 정답표}\par\vspace{4bp}
  {\small 2028학년도 대학수학능력시험 대비 모의고사 (2022 개정 교육과정)}\par\vspace{18bp}
  {\large\renewcommand{\arraystretch}{1.55}\circfont
  \begin{tabular}{|>{\centering\arraybackslash}p{14mm}|>{\centering\arraybackslash}p{20mm}|>{\centering\arraybackslash}p{14mm}||>{\centering\arraybackslash}p{14mm}|>{\centering\arraybackslash}p{20mm}|>{\centering\arraybackslash}p{14mm}|}\hline
  \textbf{번호} & \textbf{정답} & \textbf{배점} & \textbf{번호} & \textbf{정답} & \textbf{배점} \\ \hline\hline
  #1
  \end{tabular}}\par\vspace{14bp}
  배점 합계 #2점 \quad(1\textasciitilde21번 5지선다형, 22\textasciitilde30번 단답형)
  \end{center}}
